\documentclass[10pt,a4paper]{amsart}
\usepackage{amsmath,amssymb,amsthm,graphicx}
% Retained preamble for the Stacks Project chapter snapshot in this directory.
% The chapter's own preamble.tex is a copy of the shared Stacks Project preamble; the block
% below is sliced MECHANICALLY, line for line, from it. Nothing was added, removed or
% redefined; the wrapper class, the xr-hyper/hyperref wiring and the externaldocument list
% are omitted because the assembler supplies the document class.
%
% lines taken from preamble.tex: 3, 5, 9, 11, 12, 156, 157, 158, ... (51 lines total)
%
% package coverage restored from the chapter's own preamble (2026-10-13): xy, xy-2cell

\usepackage[all]{xy}
\xyoption{2cell}
\UseAllTwocells

\newenvironment{reference}{\comment}{\endcomment}
\newenvironment{slogan}{\comment}{\endcomment}
\newenvironment{history}{\comment}{\endcomment}
\theoremstyle{plain}
\newtheorem{theorem}[subsection]{Theorem}
\newtheorem{proposition}[subsection]{Proposition}
\newtheorem{lemma}[subsection]{Lemma}
\theoremstyle{definition}
\newtheorem{definition}[subsection]{Definition}
\newtheorem{example}[subsection]{Example}
\newtheorem{exercise}[subsection]{Exercise}
\newtheorem{situation}[subsection]{Situation}
\theoremstyle{remark}
\newtheorem{remark}[subsection]{Remark}
\newtheorem{remarks}[subsection]{Remarks}
\numberwithin{equation}{subsection}
\def\lim{\mathop{\mathrm{lim}}\nolimits}
\def\colim{\mathop{\mathrm{colim}}\nolimits}
\def\Spec{\mathop{\mathrm{Spec}}}
\def\Hom{\mathop{\mathrm{Hom}}\nolimits}
\def\Ext{\mathop{\mathrm{Ext}}\nolimits}
\def\SheafHom{\mathop{\mathcal{H}\!\mathit{om}}\nolimits}
\def\SheafExt{\mathop{\mathcal{E}\!\mathit{xt}}\nolimits}
\def\Sch{\mathit{Sch}}
\def\Mor{\mathop{\mathrm{Mor}}\nolimits}
\def\Ob{\mathop{\mathrm{Ob}}\nolimits}
\def\Sh{\mathop{\mathit{Sh}}\nolimits}
\def\NL{\mathop{N\!L}\nolimits}
\def\CH{\mathop{\mathrm{CH}}\nolimits}
\def\proetale{{pro\text{-}\acute{e}tale}}
\def\etale{{\acute{e}tale}}
\def\QCoh{\mathit{QCoh}}
\def\Ker{\mathop{\mathrm{Ker}}}
\def\Im{\mathop{\mathrm{Im}}}
\def\Coker{\mathop{\mathrm{Coker}}}
\def\Coim{\mathop{\mathrm{Coim}}}
\DeclareMathSymbol{\boxtimes}{\mathbin}{AMSa}{"02}
\def\QCohstack{\mathcal{QC}\!\mathit{oh}}
\def\Cohstack{\mathcal{C}\!\mathit{oh}}
\def\Spacesstack{\mathcal{S}\!\mathit{paces}}
\def\Quotfunctor{\mathrm{Quot}}
\def\Hilbfunctor{\mathrm{Hilb}}
\def\Curvesstack{\mathcal{C}\!\mathit{urves}}
\def\Polarizedstack{\mathcal{P}\!\mathit{olarized}}
\def\Complexesstack{\mathcal{C}\!\mathit{omplexes}}
\def\Pic{\mathop{\mathrm{Pic}}\nolimits}
\def\Picardstack{\mathcal{P}\!\mathit{ic}}
\def\Picardfunctor{\mathrm{Pic}}
\def\Deformationcategory{\mathcal{D}\!\mathit{ef}}
\usepackage{cleveref}
\crefname{originallemma}{original lemma}{original lemmas}
\Crefname{originallemma}{Original lemma}{Original lemmas}
\makeatletter
\newlabel{algebraic-lemma-representable-diagonal}{{045G}{source}{}{sourcepointer-1}{}}
\newlabel{algebraic-lemma-representable-diagonal@cref}{{[originallemma][0][]045G}{[1][0][]source}{}{}{}}
\makeatother
\crefname{originallemma}{original lemma}{original lemmas}
\Crefname{originallemma}{Original lemma}{Original lemmas}
\makeatletter
\newlabel{lemma-base-change-versal}{{0CXM}{source}{}{sourcepointer-2}{}}
\newlabel{lemma-base-change-versal@cref}{{[originallemma][0][]0CXM}{[1][0][]source}{}{}{}}
\makeatother
\crefname{originallemma}{original lemma}{original lemmas}
\Crefname{originallemma}{Original lemma}{Original lemmas}
\makeatletter
\newlabel{lemma-generalization-versality}{{0G2K}{source}{}{sourcepointer-3}{}}
\newlabel{lemma-generalization-versality@cref}{{[originallemma][0][]0G2K}{[1][0][]source}{}{}{}}
\makeatother
\crefname{originallemma}{original lemma}{original lemmas}
\Crefname{originallemma}{Original lemma}{Original lemmas}
\makeatletter
\newlabel{lemma-infinite-sequence-pre}{{0G2S}{source}{}{sourcepointer-4}{}}
\newlabel{lemma-infinite-sequence-pre@cref}{{[originallemma][0][]0G2S}{[1][0][]source}{}{}{}}
\makeatother
\crefname{originallemma}{original lemma}{original lemmas}
\Crefname{originallemma}{Original lemma}{Original lemmas}
\makeatletter
\newlabel{properties-lemma-countable-dense-subset}{{0G2F}{source}{}{sourcepointer-5}{}}
\newlabel{properties-lemma-countable-dense-subset@cref}{{[originallemma][0][]0G2F}{[1][0][]source}{}{}{}}
\makeatother
\crefname{originallemma}{original lemma}{original lemmas}
\Crefname{originallemma}{Original lemma}{Original lemmas}
\makeatletter
\newlabel{properties-lemma-maximal-points}{{0G2R}{source}{}{sourcepointer-6}{}}
\newlabel{properties-lemma-maximal-points@cref}{{[originallemma][0][]0G2R}{[1][0][]source}{}{}{}}
\makeatother
\crefname{originallemma}{original lemma}{original lemmas}
\Crefname{originallemma}{Original lemma}{Original lemmas}
\makeatletter
\newlabel{remark-strong-effectiveness}{{0CXT}{source}{}{sourcepointer-7}{}}
\newlabel{remark-strong-effectiveness@cref}{{[originallemma][0][]0CXT}{[1][0][]source}{}{}{}}
\makeatother
\crefname{originallemma}{original lemma}{original lemmas}
\Crefname{originallemma}{Original lemma}{Original lemmas}
\makeatletter
\newlabel{spaces-more-morphisms-lemma-smooth-formally-smooth}{{04AM}{source}{}{sourcepointer-8}{}}
\newlabel{spaces-more-morphisms-lemma-smooth-formally-smooth@cref}{{[originallemma][0][]04AM}{[1][0][]source}{}{}{}}
\makeatother
\title{Strong effectiveness in square-zero systems implies openness of versality}
\author{The Stacks Project Authors}
\date{stacks-project:artin.tex:local-snapshot-20261012; GNU FDL 1.2 or later}
\begin{document}
\maketitle
\noindent\textit{Editorial note.} Statement span is lines 3281-3295 of sources/stacks-artin/pinned-original-source.tex; proof span is lines 3297-3384, the complete author proof (86 body lines).\par\medskip
\noindent\textit{Editorial note.} Navigation stubs: 8 cited-but-not-included labels are supplied as non-load-bearing pointers with their pinned tags where the snapshot tag index has them.\par\medskip
\noindent\textit{Editorial note.} extra\_packages is empty on purpose: the retained preamble already loads xy with options and a second declaration clashes in the assembler wrapper.\par\medskip
\noindent\textit{Editorial note.} Text is retained verbatim, including the author's own typos and placeholders.\par\medskip
\noindent\textit{Editorial note.} Queue index 7 of queue\_artin.json (pb=3812, stmt 3281-3295, proof 3297-3384).\par\medskip
\noindent\textit{Editorial note.} stacks-artin wave ar01 (2026-10-14): fresh adjudication accept per lane-s3/artin\_judgments.md; pool per Chief 10:5x direct dispatch (idle-first over S12's 01:5x note, SMC precedent); artin.tex 99 rows/67 with proof judged 46 accept / 21 reject / 5 AGP-overlap avoidance rows; patch-sd-envs contract spans scanned.\par\medskip
\noindent\textit{Allowed standard prerequisites (editorial).}\begin{itemize}
\item Explicit prerequisite boundary. The cited inputs are prerequisites; the argument in the retained span is not replaced by any cited result.
\end{itemize}\medskip
\ifdefined\hypertarget\hypertarget{sourcepointer-1}{}\fi
\noindent\textit{Source reference.} Original lemma 045G: sources/stacks-artin/pinned-original-source.tex, source lines 3384--3384. Cited by the retained statement or proof but not part of the collected spans; supplied only as a navigation pointer. The argument does not depend on its content.\par\medskip
\ifdefined\hypertarget\hypertarget{sourcepointer-2}{}\fi
\noindent\textit{Source reference.} Original lemma 0CXM: sources/stacks-artin/pinned-original-source.tex, source lines 2013--2022. Cited by the retained statement or proof but not part of the collected spans; supplied only as a navigation pointer. The argument does not depend on its content.\par\medskip
\ifdefined\hypertarget\hypertarget{sourcepointer-3}{}\fi
\noindent\textit{Source reference.} Original lemma 0G2K: sources/stacks-artin/pinned-original-source.tex, source lines 3076--3084. Cited by the retained statement or proof but not part of the collected spans; supplied only as a navigation pointer. The argument does not depend on its content.\par\medskip
\ifdefined\hypertarget\hypertarget{sourcepointer-4}{}\fi
\noindent\textit{Source reference.} Original lemma 0G2S: sources/stacks-artin/pinned-original-source.tex, source lines 3176--3207. Cited by the retained statement or proof but not part of the collected spans; supplied only as a navigation pointer. The argument does not depend on its content.\par\medskip
\ifdefined\hypertarget\hypertarget{sourcepointer-5}{}\fi
\noindent\textit{Source reference.} Original lemma 0G2F: sources/stacks-artin/pinned-original-source.tex, source lines 3384--3384. Cited by the retained statement or proof but not part of the collected spans; supplied only as a navigation pointer. The argument does not depend on its content.\par\medskip
\ifdefined\hypertarget\hypertarget{sourcepointer-6}{}\fi
\noindent\textit{Source reference.} Original lemma 0G2R: sources/stacks-artin/pinned-original-source.tex, source lines 3384--3384. Cited by the retained statement or proof but not part of the collected spans; supplied only as a navigation pointer. The argument does not depend on its content.\par\medskip
\ifdefined\hypertarget\hypertarget{sourcepointer-7}{}\fi
\noindent\textit{Source reference.} Original lemma 0CXT: sources/stacks-artin/pinned-original-source.tex, source lines 3257--3267. Cited by the retained statement or proof but not part of the collected spans; supplied only as a navigation pointer. The argument does not depend on its content.\par\medskip
\ifdefined\hypertarget\hypertarget{sourcepointer-8}{}\fi
\noindent\textit{Source reference.} Original lemma 04AM: sources/stacks-artin/pinned-original-source.tex, source lines 3384--3384. Cited by the retained statement or proof but not part of the collected spans; supplied only as a navigation pointer. The argument does not depend on its content.\par\medskip
\begin{lemma}
\label{lemma-SGE-implies-openness-versality}
Let $S$ be a locally Noetherian scheme. Let $\mathcal{X}$ be a category fibred
in groupoids over $(\Sch/S)_{fppf}$. Assume
\begin{enumerate}
\item $\Delta : \mathcal{X} \to \mathcal{X} \times \mathcal{X}$ is
representable by algebraic spaces,
\item $\mathcal{X}$ has (RS*),
\item $\mathcal{X}$ is limit preserving,
\item systems $(\xi_n)$ as in Remark \ref{remark-strong-effectiveness}
where $\Ker(R_m \to R_n)$ is an ideal of square zero for all $m \geq n$
are effective.
\end{enumerate}
Then $\mathcal{X}$ satisfies openness of versality.
\end{lemma}

\begin{proof}
Choose a scheme $U$ locally of finite type over $S$,
a finite type point $u_0$ of $U$, and an object $x$ of $\mathcal{X}$
over $U$ such that $x$ is versal at $u_0$. After shrinking
$U$ we may assume $U$ is affine and $U$ maps into an affine open
$\Spec(\Lambda)$ of $S$. Let $E \subset U$ be the set of finite type
points $u$ such that $x$ is not versal at $u$. By
Lemma \ref{lemma-generalization-versality}
if $u \in E$ then $u_0$ is not a specialization of $u$.
If openness of versality does not hold, then $u_0$ is in the closure
$\overline{E}$ of $E$. By
Properties, Lemma \ref{properties-lemma-countable-dense-subset}
we may choose a countable subset $E' \subset E$ with the same closure
as $E$. By Properties, Lemma \ref{properties-lemma-maximal-points}
we may assume there are no specializations among the points of $E'$.
Observe that $E'$ has to be (countably) infinite as $u_0$
isn't the specialization of any point of $E'$ as pointed out above.
Thus we can write $E' = \{u_1, u_2, u_3, \ldots\}$, there
are no specializations among the $u_i$, and $u_0$ is in the closure
of $E'$.

\medskip\noindent
Choose $x \to x_1 \to x_2 \to \ldots$ lying over
$U \to U_1 \to U_2 \to \ldots$ as in Lemma \ref{lemma-infinite-sequence-pre}.
Write $U_n = \Spec(R_n)$ and $U = \Spec(R_0)$.
Set $R = \lim R_n$. Observe that $R \to R_0$ is surjective
with kernel an ideal of square zero. By assumption (4)
we get $\xi$ over $\Spec(R)$ whose base change to $R_n$ is $x_n$.
By assumption (3) we get that $\xi$ comes from an object $\xi'$ over
$U' = \Spec(R')$ for some finite type $\Lambda$-subalgebra
$R' \subset R$. After increasing $R'$ we may and do assume that
$R' \to R_0$ is surjective, so that $U \subset U'$ is a first order thickening.
Thus we now have
$$
x \to x_1 \to x_2 \to \ldots \to \xi'
\text{ lying over }
U \to U_1 \to U_2 \to \ldots \to U'
$$
By assumption (1) there is an algebraic space $Z$ over $S$ representing
$$
(\Sch/U)_{fppf} \times_{x, \mathcal{X}, \xi'} (\Sch/U')_{fppf}
$$
see Algebraic Stacks, Lemma \ref{algebraic-lemma-representable-diagonal}.
By construction of $2$-fibre products, a $T$-valued point of $Z$
corresponds to a triple $(a, a', \alpha)$ consisting of morphisms
$a : T \to U$, $a' : T \to U'$ and a morphism $\alpha : a^*x \to (a')^*\xi'$.
We obtain a commutative diagram
$$
\xymatrix{
U \ar[rd] \ar[rdd] \ar[rrd] \\
& Z \ar[r]_{p'} \ar[d]^p & U' \ar[d] \\
& U \ar[r] & S
}
$$
The morphism $i : U \to Z$ comes the isomorphism $x \to \xi'|_U$.
Let $z_0 = i(u_0) \in Z$. By Lemma \ref{lemma-base-change-versal}
we see that $Z \to U'$ is smooth at $z_0$. After replacing $U$ by an
affine open neighbourhood of $u_0$, replacing $U'$ by the corresponding
open, and replacing $Z$ by the intersection of the inverse images
of these opens by $p$ and $p'$, we reach the situation where
$Z \to U'$ is smooth along $i(U)$. Note that this
also involves replacing $u_n$ by a subsequence, namely
by those indices such that $u_n$ is in the open. Moreover, condition
(3) of Lemma \ref{lemma-infinite-sequence-pre}
is clearly preserved by shrinking $U$
(all of the schemes $U$, $U_n$, $U'$ have the same underlying
topological space).
Since $U \to U'$ is a first order thickening of affine schemes,
we can choose a morphism $i' : U' \to Z$
such that $p' \circ i' = \text{id}_{U'}$ and
whose restriction to $U$ is $i$
(More on Morphisms of Spaces, Lemma
\ref{spaces-more-morphisms-lemma-smooth-formally-smooth}).
Pulling back the universal morphism
$p^*x \to (p')^*\xi'$ by $i'$ we obtain a morphism
$$
\xi' \to x
$$
lying over $p \circ i' : U' \to U$ such that the composition
$$
x \to \xi' \to x
$$
is the identity. Recall that we have $x_1 \to \xi'$ lying over
the morphism $U_1 \to U'$. Composing we get a morphism
$x_1 \to x$ whose existence contradicts condition
(3) of Lemma \ref{lemma-infinite-sequence-pre}.
This contradiction finishes the proof.
\end{proof}

\end{document}
